% GENERATED FILE. Edit canonical lessons, then run npm run generate:books.
% canonical-source-hash: fnv1a64:cf1bba0756cb5822
% canonical-lessons: JA-C123-otonashii, JA-C123-natsukashii, JA-C123-kuyashii, JA-C123-ayashii, JA-C123-kuwashii

\chapter{Well-behaved, Dear, Frustrating, Suspicious, Detailed}
\label{ch:ja-well-behaved-dear-frustrating-suspicious-detailed}
\hlchaptermodality{123}

\begin{chapteropening}
\textbf{By the end of this chapter:} \emph{I can say well-behaved, dear, frustrating, suspicious and detailed.}

Complete the last lesson of chapter 123: I can say well-behaved, dear, frustrating, suspicious and detailed.
\end{chapteropening}

\section[\texorpdfstring{\ja{おとなしい} (otonashii)}{おとなしい (otonashii)}]{\ja{おとなしい} (otonashii) — well-behaved, quiet}

\label{lesson:JA-C123-otonashii}



\begin{quote}
Before the new one: say the Japanese for simple, then the Japanese for complicated.
\end{quote}

\subsection*{You'll want to know: \ja{おとなしい}}
\textbf{\ja{おとなしい}} — \emph{otonashii} — \textquotedblleft{}well-behaved, quiet\textquotedblright{}.

\begin{grammarlens}[title={What you've built}]
One more word for this chapter.
\end{grammarlens}

\subsection*{Your turn}
\begin{itemize}
\raggedright
  \item \emph{Say it:} \emph{otonashii}
  \item \emph{Say it:} \emph{otonashii}, once more
  \item \emph{Say it:} say \emph{otonashii}
  \item \emph{Recall:} say the Japanese for simple, then the Japanese for complicated, then say \emph{otonashii} again
\end{itemize}

\begin{tcolorbox}[breakable,colback=teal!4,colframe=teal!35!black,title={Before you move on}]
What does \ja{おとなしい} mean? (\textquotedblleft{}Well-behaved, quiet\textquotedblright{}.) Say it once more.
\end{tcolorbox}
\section[\texorpdfstring{\ja{なつかしい} (natsukashii)}{なつかしい (natsukashii)}]{\ja{なつかしい} (natsukashii) — dear, missed}

\label{lesson:JA-C123-natsukashii}



\begin{quote}
Before the new one: say the Japanese for complicated, then the Japanese for well-behaved.
\end{quote}

\subsection*{You'll want to know: \ja{なつかしい}}
\textbf{\ja{なつかしい}} — \emph{natsukashii} — \textquotedblleft{}dear, missed\textquotedblright{}.

\begin{grammarlens}[title={What you've built}]
One more word for this chapter.
\end{grammarlens}

\subsection*{Your turn}
\begin{itemize}
\raggedright
  \item \emph{Say it:} \emph{natsukashii}
  \item \emph{Say it:} \emph{natsukashii}, once more
  \item \emph{Say it:} say \emph{natsukashii}
  \item \emph{Recall:} say the Japanese for complicated, then the Japanese for well-behaved, then say \emph{natsukashii} again
\end{itemize}

\begin{tcolorbox}[breakable,colback=teal!4,colframe=teal!35!black,title={Before you move on}]
What does \ja{なつかしい} mean? (\textquotedblleft{}Dear, missed\textquotedblright{}.) Say it once more.
\end{tcolorbox}
\section[\texorpdfstring{\ja{くやしい} (kuyashii)}{くやしい (kuyashii)}]{\ja{くやしい} (kuyashii) — frustrating}

\label{lesson:JA-C123-kuyashii}



\begin{quote}
Before the new one: say the Japanese for well-behaved, then the Japanese for dear.
\end{quote}

\subsection*{You'll want to know: \ja{くやしい}}
\textbf{\ja{くやしい}} — \emph{kuyashii} — \textquotedblleft{}frustrating\textquotedblright{}.

\begin{grammarlens}[title={What you've built}]
One more word for this chapter.
\end{grammarlens}

\subsection*{Your turn}
\begin{itemize}
\raggedright
  \item \emph{Say it:} \emph{kuyashii}
  \item \emph{Say it:} \emph{kuyashii}, once more
  \item \emph{Say it:} say \emph{kuyashii}
  \item \emph{Recall:} say the Japanese for well-behaved, then the Japanese for dear, then say \emph{kuyashii} again
\end{itemize}

\begin{tcolorbox}[breakable,colback=teal!4,colframe=teal!35!black,title={Before you move on}]
What does \ja{くやしい} mean? (\textquotedblleft{}Frustrating\textquotedblright{}.) Say it once more.
\end{tcolorbox}
\section[\texorpdfstring{\ja{あやしい} (ayashii)}{あやしい (ayashii)}]{\ja{あやしい} (ayashii) — suspicious}

\label{lesson:JA-C123-ayashii}



\begin{quote}
Before the new one: say the Japanese for dear, then the Japanese for frustrating.
\end{quote}

\subsection*{You'll want to know: \ja{あやしい}}
\textbf{\ja{あやしい}} — \emph{ayashii} — \textquotedblleft{}suspicious\textquotedblright{}.

\begin{grammarlens}[title={What you've built}]
One more word for this chapter.
\end{grammarlens}

\subsection*{Your turn}
\begin{itemize}
\raggedright
  \item \emph{Say it:} \emph{ayashii}
  \item \emph{Say it:} \emph{ayashii}, once more
  \item \emph{Say it:} say \emph{ayashii}
  \item \emph{Recall:} say the Japanese for dear, then the Japanese for frustrating, then say \emph{ayashii} again
\end{itemize}

\begin{tcolorbox}[breakable,colback=teal!4,colframe=teal!35!black,title={Before you move on}]
What does \ja{あやしい} mean? (\textquotedblleft{}Suspicious\textquotedblright{}.) Say it once more.
\end{tcolorbox}
\section[\texorpdfstring{\ja{くわしい} (kuwashii)}{くわしい (kuwashii)}]{\ja{くわしい} (kuwashii) — detailed}

\label{lesson:JA-C123-kuwashii}



\begin{quote}
Before the new one: say the Japanese for frustrating, then the Japanese for suspicious.
\end{quote}

\subsection*{You'll want to know: \ja{くわしい}}
\textbf{\ja{くわしい}} — \emph{kuwashii} — \textquotedblleft{}detailed\textquotedblright{}.

\begin{grammarlens}[title={What you've built}]
One more word for this chapter.
\end{grammarlens}

\subsection*{Your turn}
\begin{itemize}
\raggedright
  \item \emph{Say it:} \emph{kuwashii}
  \item \emph{Say it:} \emph{kuwashii}, once more
  \item \emph{Say it:} say \emph{kuwashii}
  \item \emph{Recall:} say the Japanese for frustrating, then the Japanese for suspicious, then say \emph{kuwashii} again
\end{itemize}

\begin{tcolorbox}[breakable,colback=teal!4,colframe=teal!35!black,title={Before you move on}]
What does \ja{くわしい} mean? (\textquotedblleft{}Detailed\textquotedblright{}.) Say it once more.
\end{tcolorbox}
